%!TEX root = ../../main.tex
%% ===============================================================
%% 6.4 관측 창 구성과 정규화
%% 파일: chapters/06_features/04_window_normalization.tex
%% 상위: chapters/06_features/chapter.tex
%% 정의 라벨: sec:feat-window
%% 참조 라벨: asm:onsetcap, asm:stephold
%% 포함 객체: -
%% 인용: -
%% 상태: 초안 (docs/OUTLINE.md 참고)
%% ===============================================================

\section{관측 창 구성과 정규화}
\label{sec:feat-window}

구간 $\ell \in \{0, \dots, 5\}$의 경계는 $b_0 = t_e - 30{,}000 + 5{,}000\,\ell$, $b_1 = b_0 + 5{,}000$이고 중점은 $q_\ell = b_0 + 2{,}500$이다. 노드·전역 특징은 가정 \ref{asm:stephold}--\ref{asm:onsetcap}에 따라 $q_\ell$ 이전의 마지막 프레임을 계단 유지하고, 사건 특징은 $[b_0, b_1)$의 계수를 취한다. 결과 텐서는 $X^{\mathrm{node}} \in \R^{6\times 10\times 76}$, $X^{\mathrm{glob}} \in \R^{6\times 26}$, $X^{\mathrm{ev}} \in \R^{6\times 44}$이다.

연속 노드 특징은 학습 분할에서 추정한 통계로 표준화하되, 범주형 열과 이진 열은 정규화에서 제외한다. 챔피언 스탯과 피해 통계는 꼬리가 긴 분포이므로 $\log(1 + x)$ 변환 후 표준화한다. 위치는 지도 크기로 나눈 $[0,1]^2$ 값을 쓰며, 그래프 모델에서는 교전 중심에 대한 상대 좌표 옵션을 지원한다. 순차 모델이 사용하는 \emph{거시 시퀀스} $S \in \R^{6 \times 95}$는 전역 특징, 노드 특징의 팀별 요약, 사건 특징, 공간 요약(25)을 이어 붙인 것이며, 지도 위치에 대한 편향 학습을 막기 위해 여기서는 위치 열을 0으로 둔다.
